\documentclass[a4paper,11pt]{article}

\usepackage{geometry}
\usepackage{booktabs}
\usepackage{longtable}
\usepackage{array}
\usepackage{fontspec}
\usepackage{hyperref}
\usepackage{xcolor}
\usepackage{enumitem}
\usepackage{listings}

\setmainfont{Noto Serif CJK SC}
\setsansfont{Noto Sans CJK SC}
\setmonofont{Noto Sans Mono CJK SC}
\XeTeXlinebreaklocale "zh"
\XeTeXlinebreakskip = 0pt plus 1pt
\tolerance=9999
\emergencystretch=3em
\sloppy

\geometry{margin=2.4cm}
\hypersetup{
  colorlinks=true,
  linkcolor=blue,
  citecolor=blue,
  urlcolor=blue
}
\setlist{nosep,leftmargin=2em}

\lstdefinestyle{paperstyle}{
  basicstyle=\ttfamily\small,
  breaklines=true,
  columns=fullflexible,
  frame=single,
  framerule=0.3pt
}
\lstset{style=paperstyle}

\title{面向跨 OS 底座的可移植 Hypervisor Core：\\以 AxVisor 解耦与多底座适配为例}
\author{FreeHypervisor 项目组}
\date{}

\begin{document}
\maketitle

\begin{abstract}
Hypervisor 已经成为云计算、轻量级隔离、机密计算、嵌入式整合和安全关键系统中的基础软件组件。然而，现有 Hypervisor 往往与某一类宿主底座紧密绑定：Xen、ACRN、Bao、Jailhouse、NOVA 和 seL4 CAmkES VMM 等系统倾向于自建或依赖专用底座；KVM 和 bhyve 等 hosted hypervisor 则深度复用 Linux 或 FreeBSD 的内核机制。这两类路线分别证明了专用底座和通用 OS 复用的价值，但都没有充分回答一个问题：同一个 Hypervisor Core 能否跨不同宿主 OS 复用？

本文以 AxVisor 为对象，提出一种面向跨 OS 底座的可移植 Hypervisor 架构。我们将原本绑定 ArceOS 运行环境的 AxVisor 拆分为 OS-independent Hypervisor Core 与 OS-specific Host Adapter，并把宿主 OS 依赖收束为一组显式的 Host Capability Interface。当前实现中，这一接口边界落实为 30 个最小验证接口，覆盖任务执行、运行时安装、宿主 CPU 与 per-CPU 初始化、控制台、时间、同步、IRQ、内存分配、地址转换以及启动入口桥接等能力。进一步地，本文把完整系统边界组织为四层：host-api、runtime contract、hv-provider-api 与 device-model/control-plane，以避免把宿主基础设施、AxVisor 运行时、硬件虚拟化执行面和管理面混为一层。

我们在 ArceOS、Asterinas 和 Linux 三类差异明显的底座上实现或规划该接口。ArceOS 用于验证从原生运行环境中解耦后的兼容性；Asterinas 用于验证新型 Linux ABI-compatible OS 上的 host adapter 与 KVM-style control plane；Linux 则作为成熟通用内核验证 Host Capability Interface 是否能够映射到真实生产级 OS 机制。本文的核心结论是：Hypervisor 的可移植性不应只理解为硬件架构移植或 guest 支持，而应包含宿主 OS 运行时可移植性；这一目标需要显式的能力契约，而不是简单的 API 包装。
\end{abstract}

\section{引言}

虚拟化已经从传统服务器整合技术演化为构建现代系统的重要基础。云平台依赖 Hypervisor 实现多租户隔离；serverless 和容器场景使用轻量级虚拟机降低共享内核风险~\cite{holmes2024severifast,goethals2022functional}；机密计算和基于 VM 的 TEE 依赖虚拟化边界构建可信执行环境；嵌入式、车载和工业系统则通过 Hypervisor 在同一硬件上整合通用 OS、实时 OS 与安全关键组件~\cite{lozano2023comprehensive,cilardo2022virtualization}。随着使用场景扩展，Hypervisor 不再只是 ``运行 guest OS 的工具''，而是调度、内存、I/O、中断、时间和管理面共同组成的系统级服务。

现有系统大致沿两条路线组织 Hypervisor 与宿主底座的关系。第一类系统自建或依赖专用底座。Xen 通过 dom0 和 toolstack 组织虚拟化系统；ACRN、Bao、Jailhouse 等面向嵌入式或静态分区场景，强调可预测性和资源隔离~\cite{li2019acrn,martins2020lightweight,martins2023shedding}；NOVA 和 seL4 CAmkES VMM 则将虚拟化建立在微内核对象模型和组件框架之上。这类系统具有清晰边界和较强控制力，但需要维护一整套专门的系统栈。

第二类系统复用通用 OS。KVM 将 Linux 本身变成 Hypervisor runtime：VM 表现为进程，vCPU 由 Linux 调度，guest memory 关联用户态地址空间，控制面通过 \texttt{/dev/kvm}、fd、\texttt{mmap} 和 ioctl 暴露给 QEMU、Firecracker、Cloud Hypervisor 和 crosvm 等用户态 VMM。bhyve 在 FreeBSD 上采取类似 hosted hypervisor 路线。这条路线的优势是复用成熟 OS 的调度、内存、文件对象、事件通知和设备框架，显著降低工程成本。但其代价是虚拟化核心与具体宿主内核机制强耦合。

本文关注第三个维度：\emph{宿主 OS 可移植性}。已有工作大量讨论 guest portability 和 hardware portability，例如支持不同 guest OS、x86/Arm/RISC-V 架构或不同硬件虚拟化扩展~\cite{s2021first,sousa2024hspv}。但对于 ``Hypervisor Core 本身能否跨不同宿主 OS 复用''，现有系统给出的答案仍不充分。KVM 证明 Linux 可以成为优秀的 Hypervisor runtime，却没有将 ``Hypervisor 需要哪些宿主 OS 能力'' 提炼为可由其他 OS 实现的能力契约。

我们的核心观察是：Hypervisor 对宿主 OS 的依赖虽然语义复杂，但可以被组织成若干明确的能力域，包括任务执行、同步、时间、中断、内存分配、地址转换、平台初始化、运行时入口以及控制平面桥接。只要把这些依赖从隐式调用改写为显式 contract，Hypervisor Core 中相当一部分 VM/vCPU、guest memory、VM exit 和设备中介逻辑就不必绑定某一个 OS。

本文以 AxVisor 为例验证这一思路。AxVisor 原本运行在 ArceOS 生态中，直接依赖 ArceOS 的 \texttt{axalloc}、\texttt{axhal}、\texttt{axtask}、\texttt{axruntime}、per-CPU 机制、guard 体系以及 hypervisor mode 初始化逻辑。我们将这些依赖拆分为 host substrate、runtime substrate 和 native hypervisor substrate，并进一步用四层接口模型定义边界。当前实现将下层 host capability 收束为 30 个接口函数，作为 Linux、Asterinas 和 ArceOS 三个底座上的最小可验证边界。

本文贡献如下：

\begin{itemize}
  \item 提出 Host Capability Model，把 Hypervisor 对宿主 OS runtime 的依赖显式化，区别于传统硬件 HAL。
  \item 将 AxVisor 解耦为 portable Hypervisor Core 与 Host Adapter，并以 30 个接口函数定义当前最小可验证边界。
  \item 提出四层接口框架：host-api、runtime contract、hv-provider-api、device-model/control-plane，明确宿主能力、运行时组织、硬件虚拟化执行面和管理面之间的职责分离。
  \item 在 ArceOS、Asterinas 和 Linux 三类底座上实现或规划适配，分析 Linux 与 Asterinas 在内存、IRQ、FDT、设备所有权等方面的语义差异。
  \item 将 KVM-compatible userspace ABI 分解为 control-plane frontend 与 AxVisor provider backend，为非 Linux OS 支持 Linux VMM workflow 提供结构化路径。
\end{itemize}

\section{背景与动机}

\subsection{Hypervisor 对宿主 OS 的依赖}

一个完整 Hypervisor 不仅需要执行 VM entry 和处理 VM exit，还需要组织一套运行时。至少包括以下能力：

\begin{itemize}
  \item 执行与调度能力：承载 vCPU loop、管理 host task、处理 vCPU wakeup 和 yield。
  \item 内存能力：分配 host page、构造 guest memory backing、完成 host virtual/physical address translation。
  \item 时间能力：读取单调时间、设置 host timer、驱动 guest timer event。
  \item 中断能力：注册 host IRQ handler、处理 external interrupt、把事件转发到 guest interrupt injection 路径。
  \item 同步能力：维护 VM/vCPU 状态机、wait queue、条件等待和跨任务唤醒。
  \item 控制平面能力：向用户态 VMM 或管理工具暴露 VM 创建、vCPU 运行、内存注册和状态访问语义。
\end{itemize}

这些能力不是传统意义上的硬件抽象。它们抽象的是宿主 OS runtime 服务：任务、页框、wait queue、timer、IRQ、文件对象、mmap 对象和用户地址空间。本文称之为 Host Capability。

\subsection{KVM 的启发与局限}

KVM 的成功说明复用通用 OS 是构建 Hypervisor 的有效工程路线。Linux 已经提供调度器、内存管理、文件描述符、\texttt{eventfd}、\texttt{poll/epoll}、设备模型接入和用户态 ABI，因此 KVM 不必从零构建完整 Hypervisor OS。

但是，KVM 的宿主依赖并没有被抽象为独立接口。KVM 直接嵌入 Linux 的 task、mm、fd、ioctl、eventfd、GUP、RCU 和调度机制中。对于 ArceOS、Asterinas 或其他研究 OS 而言，复用 KVM 往往意味着同时引入大量 Linux 内核内部机制。这与 ``复用 Hypervisor Core'' 的目标并不相同。

本文继承 KVM ``复用宿主 OS'' 的思想，但希望把它从 Linux 的具体实现中抽象出来：宿主 OS 不再是某个固定内核，而是一组显式 Host Capability 的提供者。

\subsection{Linux ABI 兼容与虚拟化生态兼容的断层}

近年来，许多新型 OS 通过 Linux ABI 兼容复用应用生态，例如 Asterinas、gVisor、Fuchsia Starnix、WSL、FreeBSD Linuxulator 和 illumos LX zones。这类工作主要解决普通 Linux 应用如何运行在非 Linux 内核上。

虚拟化应用更复杂。QEMU、Firecracker、Cloud Hypervisor、crosvm 和 lkvm 虽然运行在用户态，但它们依赖的不只是 \texttt{open/read/write/mmap/ioctl} 这些 syscall 表面形式，而是 \texttt{/dev/kvm} 背后定义的虚拟化对象语义：全局 KVM fd、VM fd、vCPU fd、\texttt{KVM\_SET\_USER\_MEMORY\_REGION}、\texttt{KVM\_RUN}、vCPU run page、\texttt{ioeventfd} 和 \texttt{irqfd} 等。

因此，一个 OS 能运行普通 Linux 程序，并不意味着它能运行依赖 KVM 的 Linux VMM。本文将这称为 Linux application compatibility 与 Linux virtualization compatibility 之间的断层。

\section{问题定义}

本文围绕两个问题展开。

\textbf{问题一：Hypervisor 的宿主依赖缺少可移植抽象。}
现有系统要么自建底座，要么深度绑定通用 OS。我们希望回答：能否将 Hypervisor 对宿主 OS 的依赖提炼为显式 Host Capability Interface，使同一个 Hypervisor Core 运行在不同 OS 底座上？

\textbf{问题二：Linux VMM workflow 需要 KVM-compatible 虚拟化语义。}
Linux ABI 兼容主要解决普通应用，而依赖 \texttt{/dev/kvm} 的 VMM 还需要 VM/vCPU/memory/run page/eventfd 等对象语义。我们希望回答：能否在 portable Hypervisor Core 之上提供 KVM-compatible frontend，使非 Linux OS 或替代 provider 支持部分 Linux VMM workflow？

问题一是本文主线，问题二是建立在问题一之上的生态兼容扩展。本文不试图替代 Linux KVM，也不声称完整实现所有 KVM ioctl、VFIO、nested virtualization、live migration 或生产级性能竞争。

\section{总体设计}

\subsection{三层运行结构}

从运行路径看，系统可以抽象为三层：

\begin{lstlisting}
Linux VMMs / management tools
QEMU / Firecracker / Cloud Hypervisor / crosvm / lkvm
        |
        | northbound: KVM-compatible userspace ABI
        v
OS-independent AxVisor Core
        |
        | southbound: Host Capability Interface
        v
ArceOS / Asterinas / Linux / other kernels
\end{lstlisting}

下层 Host Capability Interface 描述 AxVisor Core 对宿主 OS runtime 的依赖；上层 KVM-compatible ABI 面向 Linux VMM workflow；AxVisor Core 位于中间，负责 VM/vCPU 生命周期、guest memory、VM exit、寄存器状态和虚拟化语义。

\subsection{四层接口模型}

为了避免把不同职责混在同一个 adapter 中，本文进一步把接口边界组织为四层：

\begin{lstlisting}
device-model / control-plane
            |
            v
      runtime contract
            |
            v
      hv-provider-api
            |
            v
          host-api
\end{lstlisting}

\textbf{host-api} 提供宿主底座基础能力，例如内存分配、地址转换、任务、同步、时间、中断、平台探测和控制台。当前 30 个接口主要位于这一层，另有少量 runtime hook 和入口胶水。

\textbf{hv-provider-api} 表达真正的硬件虚拟化执行面，包括 VM/vCPU 创建、guest memory 注册、vCPU 运行、guest 状态访问、VM exit、虚拟中断注入和 guest 时间虚拟化。它依赖 host-api 的 building blocks，但不能被 host-api 取代。

\textbf{runtime contract} 负责 AxVisor 自身运行时组织，例如 VM/vCPU runtime 注册、vCPU loop、VM exit 分发、状态机、上下文协调和跨 vCPU 请求。

\textbf{device-model/control-plane} 负责配置、镜像装载、设备模型、管理接口以及 KVM-style frontend。QEMU/Firecracker 这类用户态 VMM 位于这一层之上或与这一层交互。

这四层划分的关键是职责分离。宿主内存分配不是 guest memory 注册；host task 不是 vCPU；host IRQ 不是 guest interrupt injection；host CSR/MSR 原语不是 guest register state；\texttt{/dev/kvm} fd/ioctl/mmap 也不是底层 hypervisor provider 本身。

\section{Host Capability Interface}

\subsection{当前 30 个最小验证接口}

当前实现将 AxVisor 对宿主底座的最小依赖收束为 30 个接口。它们不是最终接口全集，而是当前原型验证中用于划定 AxVisor Core 与 Host Adapter 边界的最小集合。

\begin{longtable}{p{0.08\linewidth}p{0.40\linewidth}p{0.42\linewidth}}
\toprule
编号 & 接口 & 语义 \\
\midrule
\endhead
1 & \texttt{KernelTaskRuntime::spawn\_task} & 通过宿主执行体承载 AxVisor runtime 或 vCPU 相关任务。 \\
2 & \texttt{install\_kernel\_task\_runtime} & 安装宿主 task runtime，使 AxVisor runtime 能调用宿主任务能力。 \\
3 & runtime 启动入口胶水 & 将宿主入口桥接到 AxVisor processor 或 boot path。 \\
4 & \texttt{HostIf::get\_host\_cpu\_num} & 查询宿主 CPU 数量。 \\
5 & \texttt{HostIf::init\_percpu} & 初始化 AxVisor 所需 per-CPU 状态。 \\
6 & \texttt{HostIf::exit} & 从 AxVisor runtime 请求宿主退出或停止。 \\
7 & \texttt{ConsoleIf::write\_bytes} & 向宿主控制台输出字节流。 \\
8 & \texttt{ConsoleIf::read\_bytes} & 从宿主控制台或缓冲区读取输入。 \\
9 & \texttt{TimeIf::current\_time\_nanos} & 读取宿主单调时间。 \\
10 & \texttt{TimeIf::set\_oneshot\_timer} & 设置一次性宿主 timer，用于驱动 AxVisor timer event。 \\
11 & \texttt{SyncIf::create\_wait\_queue} & 创建宿主 wait queue 对象。 \\
12 & \texttt{SyncIf::destroy\_wait\_queue} & 销毁 wait queue 并处理生命周期。 \\
13 & \texttt{SyncIf::wait\_queue\_wait} & 当前任务在 wait queue 上等待。 \\
14 & \texttt{SyncIf::wait\_queue\_wait\_until} & 带条件或超时的等待。 \\
15 & \texttt{SyncIf::wait\_queue\_wake\_one} & 唤醒一个等待任务。 \\
16 & \texttt{SyncIf::wait\_queue\_wake\_all} & 唤醒全部等待任务。 \\
17 & \texttt{TaskIf::spawn\_task\_raw} & 创建原始宿主任务。 \\
18 & \texttt{TaskIf::join\_task} & 等待宿主任务退出并回收状态。 \\
19 & \texttt{TaskIf::current\_task} & 获取当前宿主任务身份，用于 current vCPU context 等逻辑。 \\
20 & \texttt{TaskIf::yield\_now} & 主动让出 CPU。 \\
21 & \texttt{IrqIf::handle\_irq} & 处理宿主 IRQ 或 trap vector 并转发给 AxVisor。 \\
22 & \texttt{IrqIf::register\_irq\_handler} & 注册宿主 IRQ handler。 \\
23 & \texttt{MemoryIf::alloc\_frame} & 分配单个宿主页框。 \\
24 & \texttt{MemoryIf::alloc\_contiguous\_frames} & 分配物理连续宿主页框。 \\
25 & \texttt{MemoryIf::dealloc\_frame} & 释放单个宿主页框。 \\
26 & \texttt{MemoryIf::dealloc\_contiguous\_frames} & 释放物理连续宿主页框。 \\
27 & \texttt{MemoryIf::phys\_to\_virt} & 宿主物理地址到宿主虚拟地址转换。 \\
28 & \texttt{MemoryIf::virt\_to\_phys} & 宿主虚拟地址到宿主物理地址转换。 \\
29 & runtime hook 安装链 & 安装 processor/runtime/timer/IRQ 相关 hook。 \\
30 & runtime 入口桥接链 & 将 Linux/Asterinas/ArceOS 入口接入真实 \texttt{axvisor\_core::boot::run()}。 \\
\bottomrule
\end{longtable}

这组接口的意义不在于数量本身，而在于把原先散落在 ArceOS 内部的隐式依赖转化为显式 contract。它刻画的是 AxVisor Core 当前启动、运行和处理基础事件所需的宿主能力，而不是完整虚拟化 provider API。

\subsection{接口状态}

以 Linux adapter 为例，当前 30 个接口可以分为三类。

\textbf{已经具备 Linux 语义封装的接口。}
task/sync 相关接口已经映射到 Linux \texttt{kthread}、task registry、\texttt{CondVar}、\texttt{Mutex} 和完成通知；host/console/time 相关接口已经接入 CPU 查询、per-CPU 初始化、退出入口、控制台缓冲、单调时间和 RISC-V hrtimer backend；memory 相关接口已经接到 Linux 页分配、连续页分配和地址转换记录，不是简单 stub。

\textbf{已有第一版但需要补强的接口。}
runtime hook 和入口桥接链已经存在，并且 \texttt{run()} 路径能够通过 \texttt{core\_link::boot}、\texttt{boot\_vendor\_bridge}、\texttt{vendor::axvisor\_core::boot} 和 \texttt{axvisor\_linux\_bridge\_boot\_run()} 进入真实 \texttt{axvisor\_core::boot::run()}。但是，这只能证明入口桥接成立，不能单独证明 Linux guest 启动协议完整匹配。IRQ 注册与分发也已经有骨架和 external IRQ event pipeline，但真实 IRQ 源号和 guest injection 语义仍需补强。

\textbf{必须等真实 core 和真实 workload 接入后才能验证的接口。}
\texttt{TimeIf::set\_oneshot\_timer} 已有 host timer backend，并能进入 \texttt{axvisor\_core::vmm::timer::check\_events()} 路径；但 timer event 是否满足 AxVisor Core 的真实语义，必须在真实 guest workload 下验证。类似地，Linux adapter 已经能识别、缓存和转发 external IRQ 事件，但是否能正确完成 guest interrupt injection，取决于 IRQ 源号、vCPU context 和 PLIC/APIC completion 语义。

\section{实现}

\subsection{从 ArceOS 内嵌实现到 Portable AxVisor Core}

AxVisor 原本依赖 ArceOS 的多个内部组件，包括 \texttt{axalloc}、\texttt{axhal}、\texttt{axtask}、\texttt{axruntime}、\texttt{ax-kspin} guard 体系、per-CPU 区、SBI firmware 对接、基础平台 I/O 以及 hypervisor mode 初始化逻辑。按真实运行内容看，这些依赖可以归为三类 substrate：

\begin{itemize}
  \item host substrate：内存、时间、中断、架构寄存器、trap、地址空间、控制台和平台 I/O。
  \item runtime substrate：任务、guard、per-CPU、bootstrap、wait queue 和早期初始化顺序。
  \item native hypervisor substrate：VM/vCPU、world switch、G-stage/EPT/NPT、VM exit、guest CSR/MSR 状态、虚拟中断和虚拟 timer。
\end{itemize}

解耦后的 AxVisor Core 不再直接假设 ArceOS 的具体对象模型，而是通过 Host Adapter 获取宿主能力。ArceOS adapter 用于保持原生路径；Asterinas adapter 将 OSTD 和内核能力整理成 AxVisor 可用的 facade；Linux adapter 则通过 C/Rust bridge 将 Linux kernel objects 映射到同一接口。

\subsection{ArceOS 适配}

ArceOS 是 AxVisor 的原始运行环境。其价值在于验证：解耦之后的 AxVisor Core 仍能在原生底座上运行，而不是因为抽象化破坏既有语义。ArceOS 的 unikernel-style 结构使许多能力天然接近 AxVisor 需求，例如线性地址映射、直接平台初始化、轻量任务和静态 VM 配置。

在论文评估中，ArceOS static mode 应作为 southbound portability 的基线：宿主启动后直接进入 AxVisor，VM 配置和 guest 镜像在启动前准备，不依赖复杂用户态控制平面。

\subsection{Asterinas 适配}

Asterinas 代表一种 Linux ABI-compatible 但内核架构不同的新型 OS。Asterinas 适配承担两类验证。第一，它需要实现 Host Capability Interface，证明 AxVisor Core 可以脱离 ArceOS 私有假设，在另一种内核对象模型中运行。第二，它可以承载 KVM-style control plane，验证 Linux ABI 兼容能否进一步扩展到虚拟化应用 workflow。

Asterinas 侧最关键的收口点不是简单补齐 trait，而是明确底座语义：哪些地址可进行 \texttt{virt\_to\_phys}，哪些内存保证物理连续，哪些映射依赖 linear mapping，task/guard/per-CPU 语义如何与 AxVisor runtime 对齐。

\subsection{Linux 适配}

Linux adapter 的目标不是替代 KVM，而是验证 Host Capability Interface 是否能映射到成熟通用内核。Linux 与 ArceOS/Asterinas 的差异使它成为强验证底座：如果接口只能在研究 OS 或受控 unikernel 环境中成立，它的通用性有限；如果同一接口可以映射到 Linux 的 kthread、wait queue、timer、page allocator、memremap/ioremap 和 device release 语义，则更能说明 Host Capability Model 有跨 OS 意义。

当前 Linux adapter 已经具备以下基础：

\begin{itemize}
  \item adapter 分层包括 \texttt{axvisor\_adapter\_main.rs}、\texttt{core\_link/}、\texttt{vendor/}、\texttt{axvisor\_core\_stub.rs} 和 \texttt{arch/riscv64.rs}。
  \item 大多数 host capability 接口已有 Linux 语义封装，尤其是 task、sync、host、console、time 和 memory。
  \item vendor staging 已经把 AxVisor 主体源码和一批依赖源码放入 Linux 树内工作区。
  \item 真实 \texttt{axvisor\_core::boot::run()} 入口桥接链已经成型，后续关键不再是补函数壳，而是打通 proc-macro、runtime crate、virtualization crate 和 \texttt{build.rs} 处理策略。
\end{itemize}

\subsection{Linux 与 Asterinas 的关键语义差异}

Linux 适配暴露出若干不能靠 API 包装掩盖的语义差异。

\textbf{地址转换。}
Asterinas 路径可以较自然地使用 \texttt{paddr\_to\_vaddr} 和反向线性映射。Linux loadable module 不能对任意 HPA 使用 \texttt{\_\_va()}。当前 Linux \texttt{phys\_to\_virt} 需要区分 adapter-owned frames、连续分配、guest reserved RAM、passthrough MMIO、host PLIC MMIO 和不支持的任意物理地址。guest RAM 可能来自 \texttt{reserved-memory/no-map} 区域，需要 \texttt{memremap(MEMREMAP\_WB)}；MMIO 和 PLIC 需要 \texttt{ioremap}。

\textbf{guest RAM backing 与 DMA 可见性。}
在 RISC-V passthrough 场景下，guest RAM 常采用 GPA=HPA 的 identical mapping，使 QEMU 或真实设备 DMA 能访问正确宿主物理地址。Linux host 需要在启动时把该区间从普通内存管理中排除，并保证 AxVisor 可通过受控映射访问。

\textbf{外部 IRQ 注入上下文。}
Asterinas 可以在处理 S external IRQ 时遍历 pending external interrupts 并调用 guest injection 路径。Linux 侧通常需要先 claim host PLIC IRQ、入队事件、唤醒或切入 vCPU 线程，再在合适的 vCPU context 中注入。若直接在非 vCPU IRQ context 中调用依赖 \texttt{current\_vcpu\_context} 的函数，可能无法完成注入。

\textbf{passthrough 设备所有权。}
Asterinas 不需要解绑 Linux platform driver。Linux 上如果要把 virtio-mmio、串口或其他 MMIO 设备交给 guest，必须先释放 host driver 所有权，否则 passthrough 不安全或不可用。

\textbf{host FDT 来源。}
Asterinas 可以由 boot path 暴露 FDT paddr。Linux module 无法可靠直接获取 boot DTB 物理地址，当前路径需要显式传入 \texttt{host\_fdt\_path} 或 \texttt{host\_fdt\_paddr}，并校验该 DTB 与 Linux host boot DTB 及 reserved-memory 节点一致。

\section{KVM-compatible Control Plane}

\subsection{为什么它不是 30 个 host 接口的一部分}

当前 30 个接口解决的是 AxVisor 如何在宿主内核运行，即 southbound host capability 问题。Firecracker、QEMU 或 Cloud Hypervisor 需要的是另一类接口：\texttt{/dev/kvm} 文件对象、KVM ioctl、vCPU fd \texttt{mmap(struct kvm\_run)}、VM/vCPU fd 生命周期、userspace guest memory pinning、\texttt{KVM\_RUN} 返回的 exit reason，以及 \texttt{ioeventfd/irqfd} 事件语义。

因此，KVM-compatible control plane 应作为 device-model/control-plane 层的一部分，映射到底层 AxVisor hv-provider-api，而不是塞进 host-api。

\subsection{KVM ABI 对象模型}

面向未修改 Firecracker 或其他 Linux VMM，需要提供如下对象模型：

\begin{lstlisting}
/dev/kvm global fd
  -> KVM_CREATE_VM returns VM fd
  -> VM fd KVM_CREATE_VCPU returns vCPU fd
  -> vCPU fd mmap exposes struct kvm_run
  -> vCPU fd KVM_RUN enters guest
\end{lstlisting}

这一层应把 KVM ABI 分解为两部分：

\begin{itemize}
  \item frontend 语义：ioctl number、fd 生命周期、\texttt{struct kvm\_run} 布局、memory slot、eventfd、run page 和用户页 pin/unpin。
  \item backend 语义：创建 AxVisor VM/vCPU、注册 guest memory、运行 vCPU、读取/设置 guest 状态、把 AxVisor exit 翻译为 KVM exit reason。
\end{itemize}

例如，\texttt{KVM\_SET\_USER\_MEMORY\_REGION} 在 frontend 中负责接收 userspace HVA、slot、GPA 和 flags；在 Linux provider 中还要 pin userspace pages；随后通过 backend 把这些 host pages 注册到 AxVisor guest address space。\texttt{KVM\_RUN} 则调用 AxVisor provider 的 \texttt{run\_vcpu}，再把 \texttt{MmioRead/MmioWrite/IoRead/IoWrite/Halt/SystemDown/FailEntry} 等 exit 翻译为 \texttt{KVM\_EXIT\_MMIO}、\texttt{KVM\_EXIT\_IO}、\texttt{KVM\_EXIT\_HLT}、\texttt{KVM\_EXIT\_SHUTDOWN} 或 \texttt{KVM\_EXIT\_FAIL\_ENTRY}。

\subsection{范围边界}

第一阶段 KVM-compatible frontend 不应承诺完整 KVM。合理目标是支持一条确定的 VMM workflow，例如 x86\_64 或 aarch64 Firecracker cold boot，并明确不支持 snapshot、migration、dirty logging、VFIO、nested virtualization 和完整 KVM capability parity。

对 Firecracker 而言，首阶段通常需要：

\begin{itemize}
  \item 全局 fd：API version、capability check、vCPU mmap size 和 VM 创建，对应 \texttt{KVM\_GET\_API\_VERSION}、\texttt{KVM\_CHECK\_EXTENSION}、\texttt{KVM\_GET\_VCPU\_MMAP\_SIZE}、\texttt{KVM\_CREATE\_VM}。
  \item x86\_64 架构初始化：CPUID/MSR 查询，vCPU CPUID/MSR/REGS/SREGS/FPU 设置，以及 LAPIC/IRQCHIP/PIT 相关 ioctl。
  \item VM fd：guest memory region 注册、vCPU 创建、\texttt{KVM\_IOEVENTFD} 和 \texttt{KVM\_IRQFD}。
  \item vCPU fd：\texttt{mmap} run page、\texttt{KVM\_SET\_SIGNAL\_MASK}、\texttt{KVM\_RUN}。
\end{itemize}

\section{评估计划}

本文原型的评估应围绕接口充分性、可移植性、语义完整性和开销展开。

\subsection{接口充分性}

需要证明当前 30 个 Host Capability 接口足以支撑 AxVisor Core 的关键路径，包括启动 core、初始化 per-CPU、创建或恢复 VM 配置、运行 vCPU loop、处理 timer event、处理 external IRQ、管理 guest memory 和输出 guest console。

同时也要证明哪些语义不应纳入这 30 个接口，例如 VM/vCPU provider 生命周期、guest memory registration、guest register state 和 KVM fd/ioctl/mmap。这些应进入 hv-provider-api 或 control-plane 层。

\subsection{代码复用率与适配成本}

应统计 AxVisor Core 在 ArceOS、Asterinas 和 Linux 间共享的代码比例，以及各 Host Adapter 的代码规模。适配成本不应只按函数数量衡量，还应记录语义复杂度，例如 Linux \texttt{phys\_to\_virt} 地址分类、guest RAM reserved range、IRQ context 转换和设备 driver release。

\subsection{多底座功能验证}

推荐设置三类验证：

\begin{itemize}
  \item ArceOS static mode：验证解耦后原生路径仍能启动 guest。
  \item Asterinas static/control mode：验证 Host Capability Interface 与 KVM-style workflow。
  \item Linux adapter mode：验证真实 \texttt{axvisor\_core::boot::run()} 入口、timer path、IRQ path、guest RAM backing 和 RISC-V 或 x86\_64 guest smoke。
\end{itemize}

\subsection{KVM-compatible ABI 覆盖}

对于 KVM-compatible frontend，应列出当前支持的 ioctl 子集、对象生命周期、run page mmap、memory slot、MMIO pending read completion、\texttt{ioeventfd} 和 \texttt{irqfd} 状态。最终验收不应是 ``接口返回成功''，而应是未修改 VMM 的失败点按启动顺序推进，直至目标 guest 输出可复现 PASS 标志。

\subsection{性能开销}

性能评估不应声称超越 KVM，而应回答抽象层成本是否可接受。可测指标包括 vCPU entry/exit latency、MMIO exit 往返、timer event latency、IRQ injection latency、guest memory registration cost、adapter 调用层开销和 guest 启动时间。

\section{相关工作}

\subsection{专用底座与 Type-I Hypervisor}

Xen、ACRN、Bao 和 Jailhouse 代表自建或专用底座路线。Xen 通过 dom0 和 toolstack 管理设备与 VM；ACRN 面向 IoT 和嵌入式场景，通过 Service VM 组织设备和管理功能~\cite{li2019acrn}；Bao 和 Jailhouse 强调静态分区和混合关键性系统~\cite{martins2020lightweight,martins2023shedding}。这类工作展示了专用 Hypervisor stack 的控制力，但通常需要维护完整底座和管理域。

\subsection{微内核与高保证虚拟化}

NOVA 和 seL4 CAmkES VMM 代表 microhypervisor 或微内核路线。它们关注 TCB 缩减、隔离性、组件化和形式化验证。但其虚拟化逻辑仍依赖特定微内核对象模型和组件运行时。本文不要求宿主 OS 采用特定微内核抽象，而是将宿主能力收束为 Hypervisor 所需的 runtime capability。

\subsection{复用通用 OS 的 Hypervisor}

KVM 和 bhyve 证明通用 OS 可以成为 Hypervisor runtime。KVM 深度复用 Linux 的调度、内存、fd、ioctl、eventfd 和用户态 ABI；bhyve 复用 FreeBSD 内核设施并通过 \texttt{/dev/vmm} 与用户态设备模型交互。本文继承 hosted hypervisor 的工程思想，但目标是把宿主依赖从 Linux/FreeBSD 内部机制中抽象出来。

\subsection{轻量级 VMM 与 KVM 生态}

Firecracker、Cloud Hypervisor、crosvm 和 QEMU 等用户态 VMM 构成 Linux KVM 生态。它们依赖 \texttt{/dev/kvm} 的对象语义，而不是仅依赖普通 syscall。本文的 KVM-compatible control plane 正是为了解耦这一层：让 KVM-style frontend 可以映射到 AxVisor provider，而不是要求直接移植 Linux KVM。

\subsection{Linux ABI 兼容系统}

Asterinas、gVisor、Fuchsia Starnix、WSL、FreeBSD Linuxulator 和 illumos LX zones 等系统从不同方向实现 Linux ABI 或 Linux-like environment。它们主要解决普通 Linux 应用兼容。本文进一步指出，虚拟化应用需要 KVM userspace ABI 的对象语义，因而 Linux virtualization compatibility 是比 syscall compatibility 更深的一层问题。

\subsection{硬件虚拟化与 RISC-V}

RISC-V H-extension、Arm virtualization extension 和 x86 VMX/SVM 相关工作主要关注硬件虚拟化机制和架构可移植性~\cite{s2021first,sousa2024hspv}。本文与之互补：硬件支持是 Hypervisor 的必要条件，但并不解决 Hypervisor Core 如何跨宿主 OS 复用的问题。

\section{讨论与限制}

本文工作有明确边界。第一，我们不试图在 Linux 上替代 KVM。AxVisor on Linux 的作用是验证 Host Capability Interface 在成熟 OS 上可实现，而不是竞争生产级 Linux 虚拟化栈。第二，当前 30 个接口是最小验证边界，不是最终完整接口全集。随着 KVM frontend、x86\_64 backend、aarch64 backend、VFIO 或 migration 等能力加入，hv-provider-api 和 control-plane 层会继续扩展。第三，接口存在不等于语义成立。Linux 上尤其需要严格验证 guest RAM backing、\texttt{phys\_to\_virt} 地址类别、external IRQ injection context 和 passthrough device ownership。

此外，跨 OS 可移植性不是零成本。抽象层能减少重复实现 Hypervisor Core 的成本，但会把差异集中暴露在 Host Adapter 和 provider backend 中。本文主张的不是隐藏差异，而是让差异出现在正确的层次上。

\section{结论}

本文提出并实践了一种面向跨 OS 底座的可移植 Hypervisor 架构。通过将 AxVisor 拆分为 OS-independent Core 与 OS-specific Adapter，并用 30 个 Host Capability 接口定义当前最小验证边界，我们把原本隐式绑定 ArceOS 的运行时依赖转化为可由 ArceOS、Asterinas 和 Linux 分别实现的能力契约。

进一步地，本文用四层接口模型区分 host-api、runtime contract、hv-provider-api 和 device-model/control-plane，避免把宿主 OS 基础设施、AxVisor 运行时、硬件虚拟化执行面和 KVM-style 管理面混为一体。该结构说明：Hypervisor 可移植性不仅是硬件架构适配，也包括宿主 OS runtime 的可移植性；KVM 生态兼容也不仅是 syscall ABI，而是 \texttt{/dev/kvm} 对象语义与 provider backend 的组合。

当前实现已经证明大多数 host capability 可以在 Linux 上落到真实内核对象，并且 runtime 入口桥接已经进入真实 \texttt{axvisor\_core::boot::run()} 路径。后续工作应继续完成真实 guest 启动闭环、严格化 Linux memory contract、修正 external IRQ 注入上下文，并推进 KVM-compatible frontend 到未修改 VMM 的可复现启动验证。

\bibliographystyle{plain}
\bibliography{paper_refs}

\end{document}
